% Agents décisionnels : apprentissage par renforcement (section « Apprentissage par renforcement », chapitre 3)
\subsection{Formulation de l'apprentissage}\label{sec:rl}

L'agent décide une fois par jour, au début de la journée, à partir de l'état observable de la flotte. Un épisode couvre une année, soit 365 décisions, et chaque épisode tire une nouvelle graine du simulateur : l'agent ne revoit jamais deux fois la même suite de pannes. La récompense d'une journée est l'opposé de son coût total, qui regroupe les huit postes du cahier des charges (stockage, achats, urgences, main-d'œuvre, surcoût correctif, arrêts, rebut et retouche), divisé par 150 pour ramener les récompenses vers l'unité. Maximiser la somme des récompenses revient donc exactement à minimiser le coût annuel, sans pondération arbitraire entre coût et disponibilité : les arrêts sont valorisés par la marge perdue.

L'observation ne contient que des grandeurs connues d'un responsable de maintenance réel : stock, commandes en cours et machines en attente pour chaque référence, avancement de chaque composant par rapport à son intervalle de maintenance (moyenne, maximum et part des composants au-delà de 80 et 100~\%), part des machines en service, en attente de pièce ou de technicien, techniciens libres, taux de défauts et alarmes des dernières 24 heures, avancement de l'année. Deux grandeurs d'atelier complètent cette base : le risque de panne estimé par le réseau de survie pour chaque composant, qui déclenche la maintenance de P3, et la dérive normalisée du signal de chaque composant, qui déclenche la maintenance conditionnelle de P2. Les durées de vie réelles des composants restent cachées. Toutes les grandeurs sont normalisées par des valeurs de référence du scénario, ce qui donne 55 variables, contre 39 sans ces deux blocs.

\subsection{Trois niveaux d'action}

Le cahier des charges prévoyait deux niveaux d'action (§\,5.3). Un troisième, plus simple, a été ajouté pour pouvoir appliquer un DQN classique, qui ne sait traiter qu'un petit nombre d'actions discrètes.

\begin{tabularx}{\linewidth}{@{}P{2.3cm} L P{3.6cm}@{}}
\caption{Niveaux d'action proposés aux agents}\label{tab:niveaux}\\
\toprule
\hd{Niveau} & \hd{Décision quotidienne} & \hd{Taille de l'espace} \\
\midrule
Réduit & Un seuil de maintenance commun à tous les composants et un multiplicateur de stock commun à toutes les références, chacun parmi six valeurs & 36 actions \\
N1 & Un seuil de maintenance par composant (0,6 à 1,5 fois l'intervalle, ou jamais) et un niveau de recomplètement par référence (0,5 à 3 fois le niveau de référence) & $6^8 \approx 1{,}7$ million de combinaisons \\
N2 & Quantité commandée et commande urgente pour chaque référence, et pour chaque machine : ne rien faire ou remplacer l'un de ses quatre composants & plus de $10^{76}$ combinaisons \\
\bottomrule
\end{tabularx}

Au niveau N1, l'agent règle chaque jour les paramètres d'une politique à seuils, que le simulateur applique ensuite. Ce niveau reste lisible pour un responsable de maintenance : il voit à tout moment quels seuils l'agent a retenus. Au niveau N2, l'agent prend lui-même chaque décision. Les actions impossibles ou absurdes y sont masquées : maintenance d'une machine déjà arrêtée, d'un composant dont la pièce manque ou qui n'a pas atteint la moitié de son intervalle, ou alors qu'aucun technicien n'est libre. Pour ce niveau, l'observation est complétée, pour chaque machine, par son état, le risque, la dérive et l'avancement de ses quatre composants, soit 1\,355 variables.

\subsection{Algorithmes}

\intertitre{MaskablePPO pour P4 et P5}
PPO \cite{ppo} est un algorithme acteur-critique robuste aux espaces d'actions discrets multiples. Sa variante MaskablePPO de la bibliothèque sb3-contrib \cite{sb3} exclut les actions invalides de la distribution de l'acteur, ce qui accélère nettement l'apprentissage par rapport à une pénalité \cite{masking}. Les deux agents partagent leurs réglages : huit environnements en parallèle, 365 pas par environnement entre deux mises à jour, lots de 730 transitions, taux d'apprentissage $3\times10^{-4}$, facteur d'actualisation 0,995 (horizon effectif de 200 jours, du même ordre que les délais et les durées de vie). P4 dispose de 300\,000 pas, soit environ 820 années simulées, et d'un coefficient d'entropie de 0,01 ; P5 de 300\,000 pas également, mais sans bonus d'entropie. Chaque agent est en outre entraîné sur trois graines, et l'horizon de ses épisodes croît de 90 jours à une année complète au fil de l'apprentissage ; la politique retenue est celle dont le coût annuel moyen sur cinq graines de validation est le plus faible. Ces choix, ainsi que la normalisation des récompenses et la conservation du meilleur modèle, sont justifiés plus bas. Au niveau N2, la loi masquée de sb3-contrib traite les 108 dimensions de l'action une à une et coûtait 77~ms par décision, quinze fois le temps de simulation d'une journée ; une version vectorisée, dont un test vérifie qu'elle donne les mêmes probabilités, ramène ce temps à 1,5~ms.

\intertitre{DQN au niveau réduit}
Le DQN \cite{dqn} de Stable-Baselines3 sert de référence : réseau de deux couches de 64 neurones, mémoire de 200\,000 transitions, exploration $\varepsilon$ décroissante sur les 10 premiers pourcents de l'apprentissage, 300\,000 pas.

\intertitre{Branching Dueling Q-Network}
Un DQN classique estime une valeur par action ; au niveau N2, il devrait en estimer plus de $10^{76}$. L'architecture à branches de Tavakoli et al. \cite{bdq} contourne cette explosion : un tronc commun de deux couches de 256 neurones alimente une valeur d'état et, pour chaque dimension de l'action, une tête d'avantages. Le nombre de sorties croît alors avec la somme des tailles des dimensions (près de 600) et non avec leur produit. La cible d'apprentissage suit le double DQN, en moyenne sur les branches, et les actions invalides sont masquées dans chaque branche. Ce réseau, absent des bibliothèques courantes, a été écrit directement avec PyTorch ; les sorties de toutes les branches sont rangées dans un même tableau pour éviter une boucle sur la centaine de branches. L'apprentissage dure 200\,000 pas, avec une mise à jour tous les quatre pas sur des lots de 64 transitions, une mémoire de 100\,000 transitions, un taux d'apprentissage de $10^{-4}$ et une exploration $\varepsilon$ qui décroît de 1 à 0,05 sur les 30 premiers pourcents.

\intertitre{Point de départ des agents N2}
Un agent non entraîné choisit presque uniformément entre les options de chaque dimension : au niveau N2, il commande chaque jour des quantités aléatoires et envoie en maintenance quatre machines sur cinq. Ces premières années coûtent plus de 600\,000 UR, dix fois le coût normal, et la multitude d'interventions ralentit la simulation au point de ramener l'apprentissage à une vingtaine de pas par seconde. Nous ajoutons donc, avant apprentissage, un biais aux sorties des deux agents N2 qui rend l'option « ne rien faire » largement majoritaire ; l'exploration du BDQ tire ses actions aléatoires selon ce même biais. L'agent part ainsi d'une politique prudente et apprend quand s'en écarter, ce qui ne restreint aucune décision. Le bonus d'entropie habituel a en revanche été écarté au niveau N2 : sommé sur 108 dimensions, il pousse l'acteur vers des choix uniformes, c'est-à-dire précisément vers les années absurdes que le biais initial écarte.

\intertitre{Normalisation des récompenses}
Sans normalisation, les deux agents piétinaient. Le premier apprentissage de P4, mené sans précaution particulière, en donne la mesure : le coût moyen des épisodes est passé de 53\,800 à 53\,400 UR en 300\,000 pas, autant dire qu'il n'a pas bougé. Dans les deux cas, la variance des retours expliquée par le critique restait nulle, si bien que les avantages estimés n'étaient guère que du bruit. Les retours, de l'ordre de $-200$, varient beaucoup d'une année à l'autre ; les récompenses sont donc divisées par l'écart-type courant des retours, ce qui ramène les cibles du critique vers l'unité. La variance expliquée est alors passée à plus de 0,5 dès les premières mises à jour. Les deux agents ont été réentraînés ainsi, et le modèle conservé est celui dont le coût annuel moyen des derniers épisodes est le plus faible, et non le dernier. Chaque épisode est d'abord ramené à une année, en divisant sa récompense par sa durée, sinon le curriculum retiendrait le tout premier tirage : ses épisodes de quatre-vingt-dix jours, plus courts, paraissent artificiellement moins coûteux que des années entières.

\subsection{Courbes d'apprentissage}

La figure~\ref{fig:apprentissage} suit le coût annuel moyen des derniers épisodes au fil de l'apprentissage. Ces épisodes sont joués avec la politique stochastique de l'agent, ou avec son exploration pour les deux DQN, et sur des graines chaque fois nouvelles : la courbe mesure la progression, pas la performance finale, que le chapitre~\ref{sec:evaluation} évalue sur des années communes à toutes les politiques.

\begin{figure}[H]
\centering
\pgfplotsset{courbe/.style={rapport, width=\linewidth, height=6cm,
  xlabel={Pas d'apprentissage (milliers)},
  x filter/.expression={x/1000}, y filter/.expression={y/1000}, xmin=0,
  xticklabel style={/pgf/number format/use comma, /pgf/number format/fixed, /pgf/number format/precision=0},
  yticklabel style={/pgf/number format/use comma, /pgf/number format/fixed, /pgf/number format/precision=0},
  legend style={at={(1,1)}, anchor=north east, fill=white, fill opacity=0.9, text opacity=1},
  legend cell align=left}}
\begin{minipage}[t]{0.49\linewidth}
\begin{tikzpicture}
\begin{axis}[courbe, xmax=300, ymin=50, ymax=64, ylabel={Coût annuel moyen (milliers d'UR)},
  title={\sffamily\footnotesize Niveaux réduit et N1}]
\addplot[primaire, thick, mark=none] table[x=pas, y=cout_annuel, col sep=comma] {data/courbe_maskable_ppo_n1_s1_nominal_s0.csv};
\addplot[primaire!45, mark=none, densely dashed] table[x=pas, y=cout_annuel, col sep=comma] {data/courbe_maskable_ppo_n1_sans_normalisation.csv};
\addplot[gris, mark=none] table[x=pas, y=cout_annuel, col sep=comma] {data/courbe_dqn_reduced_s1_nominal_s0.csv};
\addplot[gris!70, mark=none, densely dotted, domain=0:300000] {52136};
\legend{P4, P4 sans normalisation, DQN, P1 (évaluation)}
\end{axis}
\end{tikzpicture}
\end{minipage}\hfill
\begin{minipage}[t]{0.49\linewidth}
\begin{tikzpicture}
\begin{axis}[courbe, xmax=300, ymin=40, ymax=280,
  title={\sffamily\footnotesize Niveau N2}]
\addplot[accent, thick, mark=none] table[x=pas, y=cout_annuel, col sep=comma] {data/courbe_maskable_ppo_n2_s1_nominal_s0.csv};
\addplot[accent!60, mark=none, densely dashed] table[x=pas, y=cout_moy, col sep=comma] {data/courbe_bdq.csv};
\addplot[gris!70, mark=none, densely dotted, domain=0:300000] {52136};
\legend{P5, BDQ, P1 (évaluation)}
\end{axis}
\end{tikzpicture}
\end{minipage}
\caption{Coût annuel moyen des épisodes d'apprentissage : moyenne glissante des 100 derniers épisodes pour PPO et DQN, des 20 derniers pour le BDQ. Le pointillé rappelle le coût de P1 mesuré lors de l'évaluation finale.}
\label{fig:apprentissage}
\end{figure}

Une fois les récompenses normalisées, P5 progresse nettement : sa courbe part de 109\,200 UR et descend jusqu'à environ 66\,700 UR en fin d'apprentissage ; la durée d'apprentissage n'est donc plus ce qui le limite. P4, dont l'espace d'action est bien plus étroit, part de 62\,300 UR, culmine à 63\,500 UR vers 26\,000 pas, puis redescend et se stabilise un peu au-dessus de 53\,200 UR, sans que la normalisation change grand-chose à son niveau final (53\,400 UR sans elle) : les seuils quotidiens qu'il règle, de la même famille que ceux de P1, ne lui laissent que peu de marge. Le BDQ, lui, décroche franchement : sa courbe moyenne culmine à environ 243\,000 UR vers 103\,000 pas, avant de refluer régulièrement à mesure que l'exploration s'éteint pour terminer autour de 97\,000 UR : les valeurs de ses 108 branches sont apprises sur une récompense commune à toute la flotte, et l'effet d'une décision isolée, comme l'entretien d'une machine, se perd dans le bruit des pannes des 99 autres. Le DQN, enfin, oscille sans tendance entre 55\,100 et 61\,300 UR et termine à 56\,200 UR : un seuil commun à tous les composants et un niveau de stock commun à toutes les références ne laissent guère de marge.

La figure~\ref{fig:apprentissage} suit une seule graine par agent ; entraîner trois fois le même agent donne trois courbes différentes. La figure~\ref{fig:dispersion} rassemble pour chaque agent la moyenne des trois graines et l'étendue qui les sépare.

\begin{figure}[H]
\centering
\pgfplotsset{disp/.style={rapport, width=\linewidth, height=5.0cm,
  x filter/.expression={x/1000}, y filter/.expression={y/1000}, xmin=0, xmax=305,
  xticklabel style={/pgf/number format/use comma, /pgf/number format/fixed, /pgf/number format/precision=0},
  yticklabel style={/pgf/number format/use comma, /pgf/number format/fixed, /pgf/number format/precision=0},
  legend style={at={(0.02,0.03)}, anchor=south west, fill=white, fill opacity=0.85,
    text opacity=1, draw=none, font=\sffamily\scriptsize, legend cell align=left}}}
\begin{minipage}[t]{0.49\linewidth}
\begin{tikzpicture}
\begin{axis}[disp, ymin=52, ymax=65, ylabel={Coût annuel moyen (milliers d'UR)},
  title={\sffamily\footnotesize P4, niveau N1}]
\addplot[name path=p4hi, draw=none, forget plot] table[x=pas, y=cout_max, col sep=comma] {data/courbe_p4.csv};
\addplot[name path=p4lo, draw=none, forget plot] table[x=pas, y=cout_min, col sep=comma] {data/courbe_p4.csv};
\addplot[primaire!20] fill between[of=p4hi and p4lo];
\addplot[primaire, thick, mark=none] table[x=pas, y=cout_moy, col sep=comma] {data/courbe_p4.csv};
\legend{Étendue sur trois graines, Moyenne}
\end{axis}
\end{tikzpicture}
\end{minipage}\hfill
\begin{minipage}[t]{0.49\linewidth}
\begin{tikzpicture}
\begin{axis}[disp, ymin=62, ymax=115, title={\sffamily\footnotesize P5, niveau N2}]
\addplot[name path=p5hi, draw=none, forget plot] table[x=pas, y=cout_max, col sep=comma] {data/courbe_p5.csv};
\addplot[name path=p5lo, draw=none, forget plot] table[x=pas, y=cout_min, col sep=comma] {data/courbe_p5.csv};
\addplot[accent!20] fill between[of=p5hi and p5lo];
\addplot[accent, thick, mark=none] table[x=pas, y=cout_moy, col sep=comma] {data/courbe_p5.csv};
\end{axis}
\end{tikzpicture}
\end{minipage}

\vspace{0.7em}

\begin{minipage}[t]{0.49\linewidth}
\begin{tikzpicture}
\begin{axis}[disp, ymin=54, ymax=64, xlabel={Pas d'apprentissage (milliers)},
  ylabel={Coût annuel moyen (milliers d'UR)}, title={\sffamily\footnotesize DQN, niveau réduit}]
\addplot[name path=dqnhi, draw=none, forget plot] table[x=pas, y=cout_max, col sep=comma] {data/courbe_dqn.csv};
\addplot[name path=dqnlo, draw=none, forget plot] table[x=pas, y=cout_min, col sep=comma] {data/courbe_dqn.csv};
\addplot[gris!30] fill between[of=dqnhi and dqnlo];
\addplot[gris, thick, mark=none] table[x=pas, y=cout_moy, col sep=comma] {data/courbe_dqn.csv};
\end{axis}
\end{tikzpicture}
\end{minipage}\hfill
\begin{minipage}[t]{0.49\linewidth}
\begin{tikzpicture}
\begin{axis}[disp, ymin=60, ymax=280, xlabel={Pas d'apprentissage (milliers)},
  title={\sffamily\footnotesize BDQ, niveau N2}]
\addplot[name path=bdqhi, draw=none, forget plot] table[x=pas, y=cout_max, col sep=comma] {data/courbe_bdq.csv};
\addplot[name path=bdqlo, draw=none, forget plot] table[x=pas, y=cout_min, col sep=comma] {data/courbe_bdq.csv};
\addplot[accent!18] fill between[of=bdqhi and bdqlo];
\addplot[accent!60, thick, mark=none] table[x=pas, y=cout_moy, col sep=comma] {data/courbe_bdq.csv};
\end{axis}
\end{tikzpicture}
\end{minipage}
\caption{Coût annuel moyen sur les trois graines d'entraînement de chaque agent, et étendue entre la meilleure et la pire graine au même pas. Une enveloppe étroite signale un apprentissage reproductible : en fin d'apprentissage, l'écart entre graines vaut environ 150 UR pour le DQN, 1\,000 UR pour P4, de l'ordre du millier pour P5, et 2\,200 UR pour le BDQ — un écart qui dépasse 100\,000 UR au plus fort de sa phase instable.}
\label{fig:dispersion}
\end{figure}

Les écarts de reproductibilité suivent la taille de l'espace d'action. En fin d'apprentissage, l'étendue entre les trois graines vaut environ 150 UR pour le DQN, dont l'action tient en 36 choix, et 1\,000 UR pour P4, dont l'espace dépasse pourtant le million de combinaisons : ces deux agents convergent vers des courbes presque superposées. P5, à plus de $10^{76}$ actions effectives, garde une enveloppe de l'ordre du millier d'UR. Le BDQ se détache nettement, avec 2\,200 UR entre ses graines en fin d'apprentissage, et plus de 100\,000 UR lorsque sa courbe est au plus haut : l'initialisation des poids d'un réseau à branches et l'ordre des transitions rencontrées pèsent plus lourd que pour un acteur-critique. Ce classement ne change pas la comparaison du chapitre~\ref{sec:evaluation}, où le BDQ reste loin derrière P1 quelle que soit la graine, mais il invite à lire sa courbe moyenne avec prudence.

\intertitre{Jouer l'action la plus probable ou tirer les actions}
Une politique PPO est une loi de probabilité sur les actions. L'usage veut qu'on l'évalue en jouant son mode, l'action la plus probable de chaque dimension, mais ce mode peut s'écarter beaucoup de la politique apprise. Au niveau N2, « ne rien commander » reste l'option la plus probable de chaque référence : l'agent apprend à commander certains jours seulement, et son mode ne commande jamais. Joué ainsi, P5 coûte 401\,900 UR par an sur les graines de réglage, contre 65\,000 UR en tirant ses actions. Pour chaque agent PPO, le mode de jeu est donc choisi sur les graines de réglage de P1 et P3, jamais sur celles de l'évaluation finale, et le tirage est initialisé par la graine de la simulation pour rester reproductible. Le tirage l'emporte nettement pour P5 (65\,000 UR contre 401\,900 UR pour le mode) ; pour P4, où l'écart est plus ténu, c'est en revanche le mode qui l'emporte (54\,800 UR contre 55\,600 UR) : au niveau N1, l'action la plus probable de chaque dimension reste un seuil de maintenance sensé, alors qu'au niveau N2 le mode « ne rien faire » est inutilisable.
